%% ============================================================
%% MODELO DE ATIVIDADE · PROJETO MECÂNICO
%%
%% Copie este arquivo para atividades/ (atividade-02.tex, por exemplo),
%% inclua-o no relatorio-aacc.tex com \input{atividades/atividade-02}
%% e escreva. Cada \preencher{…} é uma pendência: troque-o pelo seu
%% texto. Cada \figuraprovisoria também: troque-a pela figura de verdade.
%%
%% Uma atividade é UMA entrega sua, com começo, meio e fim. "Participei
%% do projeto X por dois anos" não é uma atividade: são várias. Divida.
%% ============================================================
\begin{atividade}{Projeto da peça Y: concepção, cálculo e protótipo}{
  tipo         = mecanico,
  modalidade   = {},   % a chave da modalidade pedida: EXT, PES, COMP, VPC, EVT, PRD, ORG, ENS, GES, CUR, OUT
  alternativa  = {},   % a modalidade alternativa, se o Colegiado não aceitar a primeira
  horas        = {},   % as horas desta atividade, com base nos registros (Parte IV)
  inicio       = {},   % AAAA-MM
  fim          = {},   % AAAA-MM, ou: em andamento
  projeto      = {},
  papel        = {},   % Responsável, Corresponsável ou Colaborador(a)
  supervisor   = {},   % quem confirma: nome e cargo (o supervisor do projeto, o gerente)
  registros    = {},   % os códigos no SOMA, com ponto e vírgula: NBL-14; NRO-PRO-003-7
  competencias = {},   % as competências das DCN que ela desenvolveu (I a VIII): I, III, V
  disciplinas  = {},   % as disciplinas do seu curso ligadas a ela: ELE075 Eletrônica; ELE067 Circuitos
}

\begin{contexto}
\preencher{Que peça ou conjunto, para que sistema, com que requisitos --- cargas, ergonomia, massa, espaço, fabricação --- e de onde veio cada um.}
\end{contexto}

\begin{contribuicao}
\preencher{O que você concebeu, modelou, calculou, fabricou e testou; o que foi de outras pessoas.}
\end{contribuicao}

\begin{desenvolvimento}
\fichatecnica{
  ferramenta    = {},   % Ferramenta CAD/CAE e versão
  materiais     = {},   % Materiais
  fabricacao    = {},   % Processo de fabricação
  pecas         = {},   % Peças e montagem — opcional
  carregamentos = {},   % Cargas e condições de uso
  normas        = {},   % Normas consideradas — opcional
  desenhos      = {},   % Desenhos técnicos (códigos) — opcional
  arquivos      = {},   % Onde estão os arquivos
}

\preencher{As alternativas de concepção e o critério de escolha. O modelo CAD. O memorial de cálculo: o dimensionamento, as tensões, os fatores de segurança, as tolerâncias e os ajustes. A análise por elementos finitos, se houve: a malha, as condições de contorno, a convergência. A fabricação e a montagem, com os parâmetros.}

\begin{figure}[htbp]
  \centering
  \figuraprovisoria{O modelo CAD: vista em perspectiva ou explodida, com as peças identificadas.}
  % \includegraphics[width=0.9\linewidth]{figuras/cad.pdf}
  \caption{Modelo CAD do conjunto Y (vista explodida).}
  \fonte{Elaborado pelo autor.}
  \label{fig:cad}
\end{figure}

\begin{figure}[htbp]
  \centering
  \figuraprovisoria{O desenho técnico, o resultado da análise ou o protótipo fabricado.}
  % \includegraphics[width=0.9\linewidth]{figuras/analise.pdf}
  \caption{Distribuição de tensões de von Mises na peça Y sob a carga de projeto.}
  \fonte{Elaborado pelo autor.}
  \label{fig:analise}
\end{figure}
\end{desenvolvimento}

\begin{resultados}
\preencher{O protótipo atendeu aos requisitos? As medições e os ensaios, os ajustes feitos e a revisão do projeto que eles pediram.}
\end{resultados}

\begin{evidencias}
% Uma linha por evidência: \evidencia{tipo}{identificador}{o que mostra}{onde conferir}
% \evidencia{Arquivo de projeto}{<projeto>/mecanica/Y\_revB.step}{O modelo CAD da revisão B}{Drive da equipe (acesso a pedido)}
% \evidencia{Arquivo controlado}{NRO-PRO-014-<PN>}{O registro de projeto (design record) da peça}{SOMA › Arquivos}
\end{evidencias}

\begin{aprendizados}
\preencher{O que o projeto real ensinou que as disciplinas de desenho, de resistência dos materiais e de elementos de máquinas não ensinaram, e o que delas você aplicou.}
\end{aprendizados}
\end{atividade}
